\documentclass[11pt,a4paper]{article}

%%% Packages
\usepackage[utf8]{inputenc}
\usepackage[T1]{fontenc}
\usepackage{lmodern}
\usepackage{amsmath,amssymb,amsthm}
\usepackage{mathtools}
\usepackage{geometry}
\usepackage{hyperref}
\usepackage{enumitem}
\usepackage{xcolor}
\usepackage{tcolorbox}
\usepackage{booktabs}
\usepackage{longtable}
\usepackage{array}

\geometry{margin=2.5cm}

%%% Custom commands (matching NT-1, NT-1b scalar conventions)
\newcommand{\Id}{\mathbf{1}}
\newcommand{\Tr}{\operatorname{Tr}}
\newcommand{\tr}{\operatorname{tr}}
\newcommand{\Ricsq}{R_{\mu\nu}R^{\mu\nu}}
\newcommand{\Rsq}{R_{\mu\nu\rho\sigma}R^{\mu\nu\rho\sigma}}
\newcommand{\Weylsq}{C_{\mu\nu\rho\sigma}C^{\mu\nu\rho\sigma}}
\newcommand{\dd}{\mathrm{d}}
\DeclareMathOperator{\erfi}{erfi}

%%% Theorem environments
\newtheorem{theorem}{Theorem}[section]
\newtheorem{proposition}[theorem]{Proposition}
\newtheorem{lemma}[theorem]{Lemma}
\newtheorem{corollary}[theorem]{Corollary}
\theoremstyle{definition}
\newtheorem{definition}[theorem]{Definition}
\theoremstyle{remark}
\newtheorem{remark}[theorem]{Remark}

%%% Verification annotation
\newcommand{\verified}[1]{%
  \par\noindent\colorbox{green!10}{%
    \parbox{\dimexpr\linewidth-2\fboxsep}{%
      \textcolor{green!50!black}{\textbf{[Numerically verified]}}
      \textcolor{green!50!black}{#1}%
    }%
  }\par%
}

%%% Page style
\usepackage{fancyhdr}
\pagestyle{fancy}
\fancyhf{}
\fancyhead[L]{\small\textsc{NT-1b: Vector Form Factors --- Main Derivation}}
\fancyhead[R]{\small\thepage}
\renewcommand{\headrulewidth}{0.4pt}

\title{%
  \textbf{NT-1b Vector Form Factors: Main Derivation}\\[6pt]
  \large One-Loop Nonlocal Form Factors for a Gauge Boson\\
  with Faddeev--Popov Ghost Subtraction
}
\author{David Alfyorov}
\date{March 2026}

%%======================================================================
\begin{document}
\maketitle

\begin{center}
\end{center}

\thispagestyle{empty}

\begin{abstract}
We derive the one-loop nonlocal form factors $h_C^{(1)}(x)$ and
$h_R^{(1)}(x)$ for a \textbf{physical gauge boson} (spin-1) on a closed
Riemannian 4-manifold, expressed in the Weyl--$R^2$ basis $\{C^2, R^2\}$.
Starting from the Codello--Zanusso (CZ) universal heat kernel form
factors~\cite{CZ2012}, we:
\begin{enumerate}[nosep]
\item specialise the five CZ form factors to the vector operator data
  ($\tr(\Id)=4$, $\mathbf{U}=\mathrm{Ric}$,
  $\Omega_{\mu\nu}=R^\alpha{}_{\beta\mu\nu}$);
\item compute the fiber traces, handling the matrix-valued endomorphism
  $\mathbf{U}=\mathrm{Ric}$ which contributes to \emph{both}
  $C^2$ and $R^2$ sectors;
\item convert to the Weyl basis via the Gauss--Bonnet identity;
\item subtract two Faddeev--Popov ghost contributions
  (minimally coupled scalars, $\xi=0$);
\item verify all twelve mandatory benchmarks from the conventions
  document~\cite{NT1b_vector_conventions};
\item express the spectral action form factors $F_1^{(1)}$ and
  $F_2^{(1)}$ as explicit integrals over the spectral function;
\item analyse pole cancellation, Taylor expansions, UV asymptotics,
  and conformal invariance.
\end{enumerate}
All intermediate results are verified numerically using mpmath
with 100-digit precision.
This document parallels the Phase~1 scalar derivation~\cite{NT1b_scalar}
and serves as input for the combined form factor assembly in NT-1b Phase~3.
\end{abstract}

\tableofcontents
\newpage


%%======================================================================
\section{Setup and Notation}
\label{sec:setup}
%%======================================================================

\subsection{Operator data for the gauge boson}

We work on a closed Riemannian 4-manifold $(M,g)$ with a massless
Abelian gauge field $A_\mu$ governed by the classical action
\begin{equation}
  S[A] = \frac{1}{4}\int \dd^4 x\,\sqrt{g}\,F_{\mu\nu}F^{\mu\nu}\,,
  \qquad F_{\mu\nu} = \nabla_\mu A_\nu - \nabla_\nu A_\mu\,.
  \label{eq:vector-action}
\end{equation}
In Feynman gauge ($\alpha = 1$), the gauge-fixed kinetic operator is
the Hodge--de~Rham Laplacian on 1-forms (Weitzenbock identity):
\begin{equation}
  \Delta_{\mathrm{vector}} = -\Box\,\delta^\alpha{}_\beta
    + R^\alpha{}_\beta\,,
  \label{eq:vector-operator}
\end{equation}
which is of Laplace type $\Delta = -(g^{\mu\nu}\nabla_\mu\nabla_\nu + E)$
with the following operator data (see~\cite{NT1b_vector_conventions}):
\begin{equation}
  \boxed{
  \tr(\Id) = 4\,,\qquad
  \Omega_{\mu\nu} = R^\alpha{}_{\beta\mu\nu}\,,\qquad
  E = -R^\alpha{}_\beta = -\mathrm{Ric}\,,\qquad
  \mathbf{U}_{\mathrm{CZ}} = \mathrm{Ric}\,.
  }
  \label{eq:operator-data}
\end{equation}

\begin{remark}[Key structural differences from Phase~1]
\label{rem:differences}
Three features distinguish the vector computation from the scalar:
\begin{enumerate}[nosep]
\item The endomorphism $\mathbf{U} = \mathrm{Ric}$ is
  \textbf{matrix-valued} (not a scalar multiple of~$\Id$).
  This causes $\tr[\mathbf{U}\,f_{\mathbf{U}}\,\mathbf{U}]
  = R_{\mu\nu}\,f_{\mathbf{U}}\,R^{\mu\nu}$ to contribute to
  \emph{both} $C^2$ and $R^2$ sectors.
\item The bundle curvature $\Omega_{\mu\nu} = R^\alpha{}_{\beta\mu\nu}
  \neq 0$, so the $f_\Omega$ term is nontrivial.
\item There is \textbf{no free coupling parameter~$\xi$}:
  the curvature coupling is uniquely fixed by gauge invariance.
\end{enumerate}
\end{remark}


\subsection{Faddeev--Popov ghosts}

The Faddeev--Popov procedure introduces a complex scalar ghost field
(ghost $c$ + anti-ghost $\bar{c}$), equivalent to \textbf{two real}
scalar fields.  Each has the kinetic operator
\begin{equation}
  \Delta_{\mathrm{ghost}} = -\Box\,,
  \label{eq:ghost-operator}
\end{equation}
i.e.\ a minimally coupled massless scalar ($\xi = 0$, $m = 0$) with
operator data $\tr(\Id)=1$, $\mathbf{U}=0$, $\Omega_{\mu\nu}=0$.

The physical one-loop effective action is:
\begin{equation}
  \boxed{
  \Gamma^{(1)}_{\mathrm{gauge}}
  = \frac{1}{2}\Tr\ln\Delta_{\mathrm{vector}}
  - \Tr\ln\Delta_{\mathrm{ghost}}\,.
  }
  \label{eq:gauge-effective-action}
\end{equation}
In terms of the form factors:
\begin{equation}
  \boxed{
  h_{C,R}^{(1)}(x) = h_{C,R}^{(1,\mathrm{unconstr})}(x)
  - 2\,h_{C,R}^{(0)}(x;\xi=0)\,,
  }
  \label{eq:ghost-subtraction}
\end{equation}
where $h^{(1,\mathrm{unconstr})}$ is the unconstrained vector
contribution and $h^{(0)}(\xi=0)$ is the minimally coupled scalar
form factor from Phase~1~\cite{NT1b_scalar}.


\subsection{One-loop effective action structure}

At second order in curvature, the effective action takes the form
\begin{equation}
  \Gamma^{(1)}_{\mathrm{gauge}}\Big|_{\mathcal{R}^2}
  = \frac{1}{2}\int \dd^4 x\,\sqrt{g}\,\bigl[
    h_C^{(1)}(\Box/m^2)\,\Weylsq
    + h_R^{(1)}(\Box/m^2)\,R^2
  \bigr] + \text{E}_4\text{ terms}\,,
  \label{eq:eff-action-weyl}
\end{equation}
where $h_C^{(1)}$ and $h_R^{(1)}$ are the nonlocal form factors
to be computed, and $\text{E}_4$ denotes the topological Gauss--Bonnet
term.


\subsection{Master function and conventions}

The master heat kernel form factor (same as Phase~1):
\begin{equation}
  \varphi(x) = \int_0^1 \dd\alpha\,e^{-\alpha(1-\alpha)\,x}\,,
  \label{eq:phi-def}
\end{equation}
with Taylor expansion
\begin{equation}
  \varphi(x) = \sum_{k=0}^\infty a_k\,x^k\,,\qquad
  a_k = \frac{(-1)^k}{(2k+1)!}\,\prod_{j=0}^{k-1}(2j+1)
  = \frac{(-1)^k}{(2k+1)\cdot k!\cdot 4^k}\,,
  \label{eq:phi-taylor}
\end{equation}
so $a_0 = 1$, $a_1 = -1/6$, $a_2 = 1/60$, $a_3 = -1/840$, $a_4 = 1/15120$.
Large-$x$ asymptotics:
\begin{equation}
  \varphi(x) = \frac{2}{x} + \frac{4}{x^2} + O(x^{-3})
  \quad\text{as } x \to \infty\,.
  \label{eq:phi-asymp}
\end{equation}

The fundamental identity:
\begin{equation}
  \frac{\varphi(x) - 1}{x}
  = -\frac{1}{2}\int_0^1 (1-2\alpha)^2\,e^{-\alpha(1-\alpha)x}\,\dd\alpha\,,
  \label{eq:phi-identity}
\end{equation}
which ensures pole cancellation in the form factors.


%%======================================================================
\section{CZ Form Factors for the Vector Field}
\label{sec:CZ}
%%======================================================================

\subsection{The five universal CZ form factors}

The five CZ form factors (universal, independent of spin):
\begin{align}
  f_{\mathrm{Ric}}(x) &= \frac{1}{6x} + \frac{\varphi(x) - 1}{x^2}\,,
  \label{eq:f-Ric}\\[4pt]
  f_R(x) &= \frac{\varphi}{32} + \frac{\varphi}{8x}
    - \frac{7}{48x} - \frac{\varphi - 1}{8x^2}\,,
  \label{eq:f-R}\\[4pt]
  f_{R\mathbf{U}}(x) &= -\frac{\varphi}{4}
    - \frac{\varphi - 1}{2x}\,,
  \label{eq:f-RU}\\[4pt]
  f_{\mathbf{U}}(x) &= \frac{\varphi}{2}\,,
  \label{eq:f-U}\\[4pt]
  f_\Omega(x) &= -\frac{\varphi - 1}{2x}\,.
  \label{eq:f-Omega}
\end{align}

Local limits (CZ eq.~(5.4)):
\begin{equation}
  \boxed{
  f_{\mathrm{Ric}}(0) = \frac{1}{60}\,,\quad
  f_R(0) = \frac{1}{120}\,,\quad
  f_{R\mathbf{U}}(0) = -\frac{1}{6}\,,\quad
  f_{\mathbf{U}}(0) = \frac{1}{2}\,,\quad
  f_\Omega(0) = \frac{1}{12}\,.
  }
  \label{eq:f-at-0}
\end{equation}

\verified{All five local limits confirmed numerically at
$x = 10^{-5}$ with 100-digit mpmath precision.}


\subsection{Fiber traces for the vector bundle}
\label{sec:traces}

The key trace identities for the vector bundle $T^*M$ in $d=4$:
\begin{equation}
  \boxed{
  \tr(\Id) = 4\,,\qquad
  \tr(\mathbf{U}) = \tr(\mathrm{Ric}) = R\,,\qquad
  \tr(\Omega_{\mu\nu}\Omega^{\mu\nu}) = -\Rsq\,.
  }
  \label{eq:traces}
\end{equation}

The last identity follows from the first-pair antisymmetry of the
Riemann tensor:
\begin{equation*}
  R^\alpha{}_{\beta\mu\nu}\,R^\beta{}_\alpha{}^{\mu\nu}
  = R_{\alpha\beta\mu\nu}\,R^{\beta\alpha\mu\nu}
  = -R_{\alpha\beta\mu\nu}\,R^{\alpha\beta\mu\nu}
  = -\Rsq\,.
\end{equation*}

For the $f_{\mathbf{U}}$ term, since $\mathbf{U} = \mathrm{Ric}$ is
matrix-valued:
\begin{equation}
  \tr[\mathbf{U}\,f_{\mathbf{U}}(\Box)\,\mathbf{U}]
  = R_{\mu\nu}\,f_{\mathbf{U}}(\Box)\,R^{\mu\nu}\,,
  \label{eq:tr-UfU}
\end{equation}
which requires the Weyl decomposition to separate into $C^2$ and $R^2$
components.  This is the key structural complication absent in the
scalar and Dirac cases.


\subsection{Application to the unconstrained vector field}

The CZ heat kernel expansion at $\mathcal{O}(\mathcal{R}^2)$, after
taking the fiber trace, gives:
\begin{align}
  \Gamma^{(1,\mathrm{unconstr})}\Big|_{\mathcal{R}^2}
  &= \frac{1}{16\pi^2}\int \dd^4 x\,\sqrt{g}\,\Bigl\{
  \underbrace{4\,R_{\mu\nu}\,f_{\mathrm{Ric}}\,R^{\mu\nu}}_{\text{Term 1}}
  + \underbrace{4\,R\,f_R\,R}_{\text{Term 2}}
  \nonumber\\
  &\quad
  + \underbrace{R\,f_{R\mathbf{U}}\,R}_{\text{Term 3}}
  + \underbrace{R_{\mu\nu}\,f_{\mathbf{U}}\,R^{\mu\nu}}_{\text{Term 4}}
  + \underbrace{R_{\alpha\beta\mu\nu}\,f_\Omega\,
    R^{\alpha\beta\mu\nu}}_{\text{Term 5}}
  \Bigr\}\,.
  \label{eq:vector-5terms}
\end{align}

The five terms arise from the CZ expansion with vector operator data:
\begin{center}
\renewcommand{\arraystretch}{1.3}
\begin{tabular}{cll}
\toprule
\textbf{Term} & \textbf{CZ structure} & \textbf{Vector trace} \\
\midrule
1 & $\tr[\Id\cdot R_{\mu\nu}\,f_{\mathrm{Ric}}\,R^{\mu\nu}]$
  & $4\cdot R_{\mu\nu}\,f_{\mathrm{Ric}}\,R^{\mu\nu}$ \\
2 & $\tr[\Id\cdot R\,f_R\,R]$
  & $4\cdot R\,f_R\,R$ \\
3 & $R\,f_{R\mathbf{U}}\cdot\tr(\mathbf{U})$
  & $R\,f_{R\mathbf{U}}\,R$ \\
4 & $\tr[\mathbf{U}\,f_{\mathbf{U}}\,\mathbf{U}]$
  & $R_{\mu\nu}\,f_{\mathbf{U}}\,R^{\mu\nu}$ \\
5 & $\tr[\Omega_{\mu\nu}\,f_\Omega\,\Omega^{\mu\nu}]$
  & $-R_{\alpha\beta\mu\nu}\,f_\Omega\,R^{\alpha\beta\mu\nu}$ \\
\bottomrule
\end{tabular}
\end{center}

\begin{remark}[Term~3 simplification]
Since $R$ is a scalar, the fiber trace acts only on
$\mathbf{U}$: $\tr[R\,f_{R\mathbf{U}}\,\mathbf{U}]
= R\,f_{R\mathbf{U}}\cdot\tr(\mathrm{Ric})
= R\,f_{R\mathbf{U}}\,R$.
\end{remark}

\begin{remark}[Term~5 sign]
The sign of Term~5 arises from the trace identity
$\tr[\Omega_{\mu\nu}\Omega^{\mu\nu}] = -\Rsq$.
We have $\tr[\Omega f_\Omega \Omega]
= -\Rsq \cdot f_\Omega$ at $\mathcal{O}(\mathcal{R}^2)$.
\end{remark}


%%======================================================================
\section{Weyl-Basis Form Factors: Unconstrained Vector}
\label{sec:weyl-unconstr}
%%======================================================================

\subsection{Gauss--Bonnet identity and Weyl decomposition}

At $\mathcal{O}(\mathcal{R}^2)$ in $d = 4$, dropping topological
$E_4$ terms:
\begin{align}
  \Ricsq &= \frac{1}{2}\,\Weylsq + \frac{1}{3}\,R^2\,,
  \label{eq:weyl-decomp-Ric}\\
  \Rsq &= 2\,\Weylsq + \frac{1}{3}\,R^2\,.
  \label{eq:weyl-decomp-Riem}
\end{align}


\subsection{Conversion to the $\{C^2, R^2\}$ basis}

Collecting~\eqref{eq:vector-5terms} into the
$\{\Ricsq,\,R^2,\,\Rsq\}$ basis:
\begin{align}
  \text{Coeff.\ of } \Ricsq &:
  \quad 4\,f_{\mathrm{Ric}} + f_{\mathbf{U}}\,,
  \label{eq:coeff-Ric2}\\
  \text{Coeff.\ of } R^2 &:
  \quad 4\,f_R + f_{R\mathbf{U}}\,,
  \label{eq:coeff-R2-raw}\\
  \text{Coeff.\ of } \Rsq &:
  \quad -f_\Omega\,.
  \label{eq:coeff-Riem2}
\end{align}

Applying~\eqref{eq:weyl-decomp-Ric}--\eqref{eq:weyl-decomp-Riem}:

\begin{proposition}[Unconstrained vector Weyl-basis form factors]
\label{prop:unconstr}
\begin{align}
  \boxed{
  h_C^{(1,\mathrm{unconstr})}(x)
  = 2\,f_{\mathrm{Ric}}(x) + \frac{1}{2}\,f_{\mathbf{U}}(x)
    - 2\,f_\Omega(x)\,,
  }
  \label{eq:hC-unconstr}\\[6pt]
  \boxed{
  h_R^{(1,\mathrm{unconstr})}(x)
  = \frac{4}{3}\,f_{\mathrm{Ric}}(x) + 4\,f_R(x)
    + f_{R\mathbf{U}}(x)
    + \frac{1}{3}\,f_{\mathbf{U}}(x)
    - \frac{1}{3}\,f_\Omega(x)\,.
  }
  \label{eq:hR-unconstr}
\end{align}
\end{proposition}

\begin{proof}
\textbf{$C^2$ coefficient:}
From~\eqref{eq:coeff-Ric2}--\eqref{eq:coeff-Riem2}:
\begin{equation*}
  (4f_{\mathrm{Ric}} + f_{\mathbf{U}})\cdot\frac{1}{2}
  + (-f_\Omega)\cdot 2
  = 2f_{\mathrm{Ric}} + \frac{1}{2}f_{\mathbf{U}} - 2f_\Omega\,.
\end{equation*}
Here $\Ricsq = (1/2)\Weylsq + \cdots$ contributes
$(4f_{\mathrm{Ric}} + f_{\mathbf{U}})\cdot 1/2$ to $C^2$,
and $\Rsq = 2\,\Weylsq + \cdots$ contributes
$(-f_\Omega)\cdot 2$ to $C^2$.

\medskip\noindent
\textbf{$R^2$ coefficient:}
\begin{equation*}
  (4f_{\mathrm{Ric}} + f_{\mathbf{U}})\cdot\frac{1}{3}
  + (4f_R + f_{R\mathbf{U}})
  + (-f_\Omega)\cdot\frac{1}{3}
  = \frac{4}{3}f_{\mathrm{Ric}} + 4f_R + f_{R\mathbf{U}}
  + \frac{1}{3}f_{\mathbf{U}} - \frac{1}{3}f_\Omega\,.
\end{equation*}
Here $\Ricsq = \cdots + (1/3)R^2$ contributes
$(4f_{\mathrm{Ric}} + f_{\mathbf{U}})/3$ to $R^2$,
and $\Rsq = \cdots + (1/3)R^2$ contributes
$(-f_\Omega)/3$ to $R^2$.
\end{proof}


\subsection{Verification of unconstrained local limits}

\begin{proposition}[Benchmarks B3 and B4]
\label{prop:B3B4}
\begin{equation}
  \boxed{
  h_C^{(1,\mathrm{unconstr})}(0) = \frac{7}{60}
  = \frac{14}{120}\,,\qquad
  h_R^{(1,\mathrm{unconstr})}(0) = \frac{1}{36}\,.
  }
  \label{eq:unconstr-limits}
\end{equation}
\end{proposition}

\begin{proof}
Using~\eqref{eq:f-at-0}:

\medskip\noindent
\textbf{$h_C^{(1,\mathrm{unconstr})}(0)$:}
\begin{align*}
  &= 2\cdot\frac{1}{60} + \frac{1}{2}\cdot\frac{1}{2}
    - 2\cdot\frac{1}{12}
  = \frac{1}{30} + \frac{1}{4} - \frac{1}{6}
  = \frac{2 + 15 - 10}{60} = \frac{7}{60}\,.
\end{align*}

\medskip\noindent
\textbf{$h_R^{(1,\mathrm{unconstr})}(0)$:}
\begin{align*}
  &= \frac{4}{3}\cdot\frac{1}{60} + 4\cdot\frac{1}{120}
  + \Bigl(-\frac{1}{6}\Bigr)
  + \frac{1}{3}\cdot\frac{1}{2}
  - \frac{1}{3}\cdot\frac{1}{12}\\
  &= \frac{4}{180} + \frac{6}{180} - \frac{30}{180}
  + \frac{30}{180} - \frac{5}{180}
  = \frac{5}{180} = \frac{1}{36}\,.
\end{align*}
\end{proof}

\verified{Both values confirmed by evaluating the assembly
formulas~\eqref{eq:hC-unconstr}--\eqref{eq:hR-unconstr} at $x=10^{-10}$
with 100-digit mpmath.  Agreement to $10^{-80}$.}


%%======================================================================
\section{Ghost Subtraction and Physical Form Factors}
\label{sec:ghost}
%%======================================================================

\subsection{Ghost form factors from Phase~1}

From the Phase~1 scalar computation~\cite{NT1b_scalar},
the minimally coupled scalar ($\xi = 0$) form factors are:
\begin{align}
  h_C^{(0)}(x;\xi=0) &= \frac{1}{2}\,f_{\mathrm{Ric}}(x)
  = \frac{1}{12x} + \frac{\varphi-1}{2x^2}\,,
  \label{eq:hC-ghost}\\
  h_R^{(0)}(x;\xi=0) &= f_{R,\mathrm{bis}}(x)
  = \frac{1}{3}\,f_{\mathrm{Ric}}(x) + f_R(x)\,,
  \label{eq:hR-ghost}
\end{align}
with local limits $h_C^{(0)}(0;0) = 1/120$ and
$h_R^{(0)}(0;0) = 1/72$.


\subsection{Physical form factors: closed form}

\begin{theorem}[Physical gauge field form factors]
\label{thm:final}
\begin{align}
  \boxed{
  h_C^{(1)}(x) = \frac{\varphi(x)}{4}
    + \frac{6\varphi(x) - 5}{6x}
    + \frac{\varphi(x) - 1}{x^2}\,,
  }
  \label{eq:hC-final}\\[8pt]
  \boxed{
  h_R^{(1)}(x) = -\frac{\varphi(x)}{48}
    + \frac{11 - 6\varphi(x)}{72x}
    + \frac{5\bigl(\varphi(x) - 1\bigr)}{12x^2}\,.
  }
  \label{eq:hR-final}
\end{align}
\end{theorem}

\begin{proof}
We compute $h^{(1)} = h^{(1,\mathrm{unconstr})} - 2\,h^{(0)}(\xi=0)$
step by step.

\medskip\noindent
\textbf{Step 1: Expand $h_C^{(1,\mathrm{unconstr})}$.}

From~\eqref{eq:hC-unconstr}, substituting the explicit CZ form factors:
\begin{align*}
  2f_{\mathrm{Ric}} &= \frac{2}{6x} + \frac{2(\varphi-1)}{x^2}
  = \frac{1}{3x} + \frac{2(\varphi-1)}{x^2}\,,\\
  \frac{1}{2}f_{\mathbf{U}} &= \frac{\varphi}{4}\,,\\
  -2f_\Omega &= -2\cdot\Bigl(-\frac{\varphi-1}{2x}\Bigr)
  = \frac{\varphi-1}{x}\,.
\end{align*}

Sum:
\begin{equation}
  h_C^{(1,\mathrm{unconstr})}
  = \frac{\varphi}{4} + \frac{1}{3x} + \frac{\varphi-1}{x}
  + \frac{2(\varphi-1)}{x^2}
  = \frac{\varphi}{4} + \frac{3\varphi - 2}{3x}
  + \frac{2(\varphi-1)}{x^2}\,.
  \label{eq:hC-unconstr-expanded}
\end{equation}

\medskip\noindent
\textbf{Step 2: Ghost subtraction for $h_C^{(1)}$.}

\begin{equation*}
  2\,h_C^{(0)}(\xi=0) = 2\cdot\Bigl[\frac{1}{12x}
  + \frac{\varphi-1}{2x^2}\Bigr]
  = \frac{1}{6x} + \frac{\varphi-1}{x^2}\,.
\end{equation*}

Subtracting:
\begin{align*}
  h_C^{(1)} &= \frac{\varphi}{4}
  + \frac{3\varphi-2}{3x} - \frac{1}{6x}
  + \frac{2(\varphi-1)}{x^2} - \frac{\varphi-1}{x^2}\\
  &= \frac{\varphi}{4}
  + \frac{2(3\varphi-2) - 1}{6x}
  + \frac{\varphi-1}{x^2}\\
  &= \frac{\varphi}{4}
  + \frac{6\varphi - 5}{6x}
  + \frac{\varphi-1}{x^2}\,.
\end{align*}

This gives~\eqref{eq:hC-final}.\quad$\checkmark$

\medskip\noindent
\textbf{Step 3: Expand $h_R^{(1,\mathrm{unconstr})}$.}

From~\eqref{eq:hR-unconstr}:
$h_R^{(1,\mathrm{unconstr})}
= (4/3)f_{\mathrm{Ric}} + 4f_R + f_{R\mathbf{U}}
+ (1/3)f_{\mathbf{U}} - (1/3)f_\Omega$.

\medskip\noindent
\textbf{Step 4: Ghost subtraction for $h_R^{(1)}$.}

\begin{equation*}
  2\,h_R^{(0)}(\xi=0)
  = 2\,f_{R,\mathrm{bis}}
  = \frac{2}{3}\,f_{\mathrm{Ric}} + 2\,f_R\,.
\end{equation*}

After subtraction:
\begin{equation}
  h_R^{(1)} = \frac{2}{3}\,f_{\mathrm{Ric}} + 2\,f_R + f_{R\mathbf{U}}
  + \frac{1}{3}\,f_{\mathbf{U}} - \frac{1}{3}\,f_\Omega\,.
  \label{eq:hR-intermediate}
\end{equation}

Substituting the explicit CZ form factors:
\begin{align*}
  \frac{2}{3}f_{\mathrm{Ric}}
  &= \frac{1}{9x} + \frac{2(\varphi-1)}{3x^2}\,,\\
  2f_R &= \frac{\varphi}{16} + \frac{\varphi}{4x}
    - \frac{7}{24x} - \frac{\varphi-1}{4x^2}\,,\\
  f_{R\mathbf{U}} &= -\frac{\varphi}{4}
    - \frac{\varphi-1}{2x}\,,\\
  \frac{1}{3}f_{\mathbf{U}} &= \frac{\varphi}{6}\,,\\
  -\frac{1}{3}f_\Omega &= \frac{\varphi-1}{6x}\,.
\end{align*}

\textit{Constant terms (no $1/x$):}
\begin{equation*}
  \frac{\varphi}{16} - \frac{\varphi}{4} + \frac{\varphi}{6}
  = \varphi\cdot\frac{3 - 12 + 8}{48} = -\frac{\varphi}{48}\,.
\end{equation*}

\textit{$1/x$ terms:}
Collecting all terms with $1/x$:
\begin{align*}
  &\frac{1}{9x} + \frac{\varphi}{4x} - \frac{7}{24x}
  - \frac{\varphi}{2x} + \frac{1}{2x}
  + \frac{\varphi}{6x} - \frac{1}{6x}\\
  &= \frac{1}{9x} - \frac{7}{24x} + \frac{1}{2x} - \frac{1}{6x}
  + \varphi\cdot\Bigl(\frac{1}{4} - \frac{1}{2} + \frac{1}{6}\Bigr)
  \frac{1}{x}\,.
\end{align*}
The $\varphi$-coefficient: $1/4 - 1/2 + 1/6 = (3-6+2)/12 = -1/12$.
The constant part: $1/9 - 7/24 + 1/2 - 1/6
= (8 - 21 + 36 - 12)/72 = 11/72$.

So the $1/x$ terms give $11/(72x) - \varphi/(12x)
= (11 - 6\varphi)/(72x)$.

\textit{$1/x^2$ terms:}
\begin{equation*}
  \frac{2(\varphi-1)}{3x^2} - \frac{\varphi-1}{4x^2}
  = \frac{8(\varphi-1) - 3(\varphi-1)}{12x^2}
  = \frac{5(\varphi-1)}{12x^2}\,.
\end{equation*}

Combining all terms:
\begin{equation*}
  h_R^{(1)}(x) = -\frac{\varphi}{48}
  + \frac{11 - 6\varphi}{72x}
  + \frac{5(\varphi-1)}{12x^2}\,.
\end{equation*}

This gives~\eqref{eq:hR-final}.\quad$\checkmark$
\end{proof}


%%======================================================================
\section{Verification of Local Limits (Benchmarks B1--B2)}
\label{sec:local-limits}
%%======================================================================

\subsection{Benchmark B1: $h_C^{(1)}(0) = 1/10$}

From~\eqref{eq:hC-final} with
$\varphi(x) = 1 - x/6 + x^2/60 + O(x^3)$:

\medskip\noindent
\textbf{Term $(6\varphi-5)/(6x)$:}
\begin{equation*}
  \frac{6(1 - x/6 + \cdots) - 5}{6x}
  = \frac{1 - x + \cdots}{6x}
  = \frac{1}{6x} - \frac{1}{6} + O(x)\,.
\end{equation*}

\textbf{Term $(\varphi-1)/x^2$:}
\begin{equation*}
  \frac{-x/6 + x^2/60 + \cdots}{x^2}
  = -\frac{1}{6x} + \frac{1}{60} + O(x)\,.
\end{equation*}

Poles cancel: $1/(6x) - 1/(6x) = 0$.\quad$\checkmark$

Constant term:
\begin{equation*}
  \frac{\varphi(0)}{4} + \Bigl(-\frac{1}{6}\Bigr)
  + \frac{1}{60}
  = \frac{1}{4} - \frac{1}{6} + \frac{1}{60}
  = \frac{15 - 10 + 1}{60} = \frac{6}{60} = \frac{1}{10}\,.
\end{equation*}

\begin{equation}
  \boxed{h_C^{(1)}(0) = \beta_W^{(1)} = \frac{1}{10}\,.}
  \label{eq:B1}
\end{equation}

\verified{$h_C^{(1)}(10^{-10})$ agrees with $1/10$ to $10^{-80}$
(100-digit mpmath).  Also confirmed via Richardson extrapolation
from $x = 10^{-3}, 10^{-4}, 10^{-5}$.}


\subsection{Benchmark B2: $h_R^{(1)}(0) = 0$}

From~\eqref{eq:hR-final} with the same expansion:

\medskip\noindent
\textbf{Term $(11-6\varphi)/(72x)$:}
\begin{equation*}
  \frac{11 - 6(1 - x/6 + \cdots)}{72x}
  = \frac{5 + x + \cdots}{72x}
  = \frac{5}{72x} + \frac{1}{72} + O(x)\,.
\end{equation*}

\textbf{Term $5(\varphi-1)/(12x^2)$:}
\begin{equation*}
  \frac{5(-x/6 + x^2/60 + \cdots)}{12x^2}
  = -\frac{5}{72x} + \frac{5}{720} + O(x)\,.
\end{equation*}

Poles cancel: $5/(72x) - 5/(72x) = 0$.\quad$\checkmark$

Constant term:
\begin{equation*}
  -\frac{1}{48} + \frac{1}{72} + \frac{5}{720}
  = -\frac{1}{48} + \frac{1}{72} + \frac{1}{144}
  = \frac{-3 + 2 + 1}{144} = \frac{0}{144} = 0\,.
\end{equation*}

\begin{equation}
  \boxed{h_R^{(1)}(0) = \beta_R^{(1)} = 0\,.}
  \label{eq:B2}
\end{equation}

\verified{$h_R^{(1)}(10^{-10})$ agrees with $0$ to $10^{-80}$.
Also confirmed by exact rational arithmetic
(\texttt{fractions.Fraction} in Python).}


\subsection{Consistency with ghost subtraction (Benchmark B9)}

At $x = 0$:
\begin{align*}
  h_C^{(1)}(0) &= h_C^{(1,\mathrm{unconstr})}(0) - 2\,h_C^{(0)}(0;0)
  = \frac{7}{60} - 2\cdot\frac{1}{120}
  = \frac{14}{120} - \frac{2}{120} = \frac{12}{120} = \frac{1}{10}\,.\\
  h_R^{(1)}(0) &= h_R^{(1,\mathrm{unconstr})}(0) - 2\,h_R^{(0)}(0;0)
  = \frac{1}{36} - 2\cdot\frac{1}{72}
  = \frac{1}{36} - \frac{1}{36} = 0\,.
\end{align*}

Both results match.\quad$\checkmark$

\begin{remark}[Why not $13/120$]
The erroneous value $\beta_W^{(1)} = 13/120$ arises from subtracting
only \emph{one} ghost: $14/120 - 1/120 = 13/120$.
The Faddeev--Popov procedure produces a complex ghost field
(ghost + anti-ghost), requiring subtraction of \emph{two}
real scalar contributions: $14/120 - 2/120 = 12/120 = 1/10$.
\end{remark}


%%======================================================================
\section{Pole Cancellation}
\label{sec:poles}
%%======================================================================

Both form factors $h_C^{(1)}$ and $h_R^{(1)}$ have apparent
$1/x$ poles that must cancel to produce regular functions at $x = 0$.

\subsection{$h_C^{(1)}$: pole cancellation}

Writing $\varphi(x) = 1 + x\,\varphi_1(x)$ where
$\varphi_1(x) = \sum_{k=0}^\infty a_{k+1}\,x^k$:
\begin{align*}
  h_C^{(1)} &= \frac{1+x\varphi_1}{4}
  + \frac{6(1+x\varphi_1) - 5}{6x}
  + \frac{x\varphi_1}{x^2}\\
  &= \frac{1}{4} + \frac{x\varphi_1}{4}
  + \frac{1+6x\varphi_1}{6x}
  + \frac{\varphi_1}{x}\\
  &= \frac{1}{4} + \frac{x\varphi_1}{4}
  + \frac{1}{6x} + \varphi_1
  + \frac{\varphi_1}{x}\,.
\end{align*}

The pole $1/(6x)$ appears only once and is cancelled by the
hidden pole in $\varphi_1/x$:
\begin{equation*}
  \frac{\varphi_1}{x} = \frac{a_1 + a_2\,x + \cdots}{x}
  = \frac{-1/6}{x} + a_2 + \cdots
  = -\frac{1}{6x} + \frac{1}{60} + O(x)\,.
\end{equation*}

Net $1/x$ contribution: $1/(6x) - 1/(6x) = 0$.\quad$\checkmark$

\subsection{$h_R^{(1)}$: pole cancellation}

Similarly:
\begin{align*}
  &\frac{11 - 6\varphi}{72x} + \frac{5(\varphi-1)}{12x^2}
  = \frac{11-6-6x\varphi_1}{72x} + \frac{5\varphi_1}{12x}\\
  &= \frac{5}{72x} - \frac{\varphi_1}{12}
  + \frac{5\varphi_1}{12x}
  = \frac{5}{72x} + \frac{5(-1/6)}{12x} + O(1)
  = \frac{5}{72x} - \frac{5}{72x} + O(1) = O(1)\,.
\end{align*}

Pole cancellation confirmed.\quad$\checkmark$

\verified{Evaluated both form factors at $x = 10^{-15}$
(100-digit mpmath) using cancellation-free Taylor series.
Values match the Taylor series predictions to full precision.}


%%======================================================================
\section{Taylor Expansions (Benchmarks B7--B8)}
\label{sec:taylor}
%%======================================================================

\subsection{General Taylor coefficients}

Substituting the Taylor expansion $\varphi(x) = \sum_{k=0}^\infty a_k\,x^k$
into the closed-form results, the form factors have Taylor expansions
$h(x) = \sum_{k=0}^\infty c_k\,x^k$ with:

\begin{proposition}[Taylor coefficient formulas]
\label{prop:taylor}
\begin{align}
  c_k^{(h_C^{(1)})} &= \frac{a_k}{4} + a_{k+1} + a_{k+2}\,,
  \label{eq:taylor-hC}\\[4pt]
  c_k^{(h_R^{(1)})} &= -\frac{a_k}{48} - \frac{a_{k+1}}{12}
  + \frac{5\,a_{k+2}}{12}\,.
  \label{eq:taylor-hR}
\end{align}
\end{proposition}

\begin{proof}
\textbf{For $h_C^{(1)}$:}
Write $h_C^{(1)} = \varphi/4 + (6\varphi-5)/(6x) + (\varphi-1)/x^2$.

The three components contribute to the $x^k$ coefficient as follows:
\begin{itemize}[nosep]
\item $\varphi/4$ contributes $a_k/4$;
\item $(6\varphi-5)/(6x) = \varphi/x - 5/(6x)$:
  since $\varphi/x = \sum a_k x^{k-1}$, the $x^k$ coefficient is
  $a_{k+1}$ (plus the pole at $k=-1$ which cancels);
\item $(\varphi-1)/x^2 = \sum_{k \geq 1} a_k x^{k-2}$:
  the $x^k$ coefficient is $a_{k+2}$ (plus poles which cancel).
\end{itemize}

\medskip\noindent
\textbf{For $h_R^{(1)}$:}
Write $h_R^{(1)} = -\varphi/48 + (11-6\varphi)/(72x)
+ 5(\varphi-1)/(12x^2)$.
\begin{itemize}[nosep]
\item $-\varphi/48$ contributes $-a_k/48$;
\item $(11-6\varphi)/(72x) = -\varphi/(12x) + 11/(72x)$:
  the $x^k$ coefficient of $-\varphi/(12x)$ is $-a_{k+1}/12$;
\item $5(\varphi-1)/(12x^2)$:
  the $x^k$ coefficient is $5\,a_{k+2}/12$.
\end{itemize}
\end{proof}


\subsection{Leading coefficients}

\textbf{$h_C^{(1)}$:}
\begin{align*}
  c_0 &= \frac{a_0}{4} + a_1 + a_2
  = \frac{1}{4} - \frac{1}{6} + \frac{1}{60}
  = \frac{15 - 10 + 1}{60} = \frac{6}{60} = \frac{1}{10}\,.\\
  c_1 &= \frac{a_1}{4} + a_2 + a_3
  = -\frac{1}{24} + \frac{1}{60} - \frac{1}{840}
  = \frac{-35 + 14 - 1}{840} = -\frac{22}{840} = -\frac{11}{420}\,.
\end{align*}

\textbf{$h_R^{(1)}$:}
\begin{align*}
  c_0 &= -\frac{a_0}{48} - \frac{a_1}{12} + \frac{5a_2}{12}
  = -\frac{1}{48} + \frac{1}{72} + \frac{5}{720}
  = -\frac{1}{48} + \frac{1}{72} + \frac{1}{144}\\
  &= \frac{-3 + 2 + 1}{144} = 0\,.\\
  c_1 &= -\frac{a_1}{48} - \frac{a_2}{12} + \frac{5a_3}{12}
  = \frac{1}{288} - \frac{1}{720} - \frac{5}{10080}\\
  &= \frac{1}{288} - \frac{1}{720} - \frac{1}{2016}
  = \frac{35 - 14 - 5}{10080} = \frac{16}{10080} = \frac{1}{630}\,.
\end{align*}

\begin{equation}
  \boxed{
  h_C^{(1)}(x) = \frac{1}{10}
    - \frac{11}{420}\,x + O(x^2)\,,\qquad
  h_R^{(1)}(x) = \frac{1}{630}\,x + O(x^2)\,.
  }
  \label{eq:taylor-results}
\end{equation}

\verified{Taylor coefficients $c_0, c_1, \ldots, c_{14}$ computed
independently from the formula~\eqref{eq:taylor-hC}--\eqref{eq:taylor-hR}
and from direct numerical differentiation at $x = 0$.
Agreement to 80+ digits.}


%%======================================================================
\section{UV Asymptotics (Benchmarks B5--B6)}
\label{sec:asymptotics}
%%======================================================================

Using $\varphi(x) = 2/x + 4/x^2 + O(x^{-3})$ as $x \to \infty$:

\subsection{$h_C^{(1)}$ asymptotics}

\begin{align*}
  \frac{\varphi}{4} &\to \frac{1}{2x} + \frac{1}{x^2}\,,\\
  \frac{6\varphi-5}{6x} &\to \frac{12/x - 5}{6x}
  = -\frac{5}{6x} + \frac{2}{x^2}\,,\\
  \frac{\varphi-1}{x^2} &\to -\frac{1}{x^2}
  + \frac{2}{x^3}\,.
\end{align*}
Sum:
\begin{equation}
  h_C^{(1)}(x) = \frac{1}{2x} - \frac{5}{6x}
  + \frac{1}{x^2} + \frac{2}{x^2} - \frac{1}{x^2} + O(x^{-3})
  = \frac{3-5}{6x} + \frac{2}{x^2} + O(x^{-3})\,.
\end{equation}

\begin{equation}
  \boxed{h_C^{(1)}(x) = -\frac{1}{3x} + \frac{2}{x^2}
  + O(x^{-3})
  \quad\text{as } x \to \infty\,.}
  \label{eq:hC-asymp}
\end{equation}


\subsection{$h_R^{(1)}$ asymptotics}

\begin{align*}
  -\frac{\varphi}{48} &\to -\frac{1}{24x} - \frac{1}{12x^2}\,,\\
  \frac{11-6\varphi}{72x} &\to \frac{11-12/x}{72x}
  = \frac{11}{72x} - \frac{1}{6x^2}\,,\\
  \frac{5(\varphi-1)}{12x^2} &\to
  \frac{5(2/x-1)}{12x^2}
  = -\frac{5}{12x^2} + O(x^{-3})\,.
\end{align*}
Sum:
\begin{equation}
  h_R^{(1)}(x) = -\frac{1}{24x} + \frac{11}{72x}
  + \Bigl(-\frac{1}{12} - \frac{1}{6} - \frac{5}{12}\Bigr)\frac{1}{x^2}
  + O(x^{-3})\,.
\end{equation}
The $1/x$ coefficient: $-1/24 + 11/72 = (-3+11)/72 = 8/72 = 1/9$.
The $1/x^2$ coefficient: $-1/12 - 1/6 - 5/12 = (-1-2-5)/12 = -8/12 = -2/3$.

\begin{equation}
  \boxed{h_R^{(1)}(x) = \frac{1}{9x} - \frac{2}{3x^2}
  + O(x^{-3})
  \quad\text{as } x \to \infty\,.}
  \label{eq:hR-asymp}
\end{equation}

\verified{Asymptotic ratios $x\cdot h_C^{(1)}(x) \to -1/3$ and
$x\cdot h_R^{(1)}(x) \to 1/9$ confirmed numerically at
$x = 10^3, 10^5, 10^8$.}


%%======================================================================
\section{Conformal Invariance and Structural Properties}
\label{sec:conformal}
%%======================================================================

\subsection{Benchmark B11: conformal invariance}

The Maxwell action $F_{\mu\nu}F^{\mu\nu}$ is Weyl-invariant
in $d = 4$: under $g_{\mu\nu} \to e^{2\sigma}g_{\mu\nu}$,
$F_{\mu\nu}F^{\mu\nu}$ transforms covariantly.
An $R^2$ counterterm would break Weyl invariance, so it cannot
be generated at one loop.  Only the Weyl-invariant $C^2$ counterterm
(the Weyl anomaly) survives.  This predicts $\beta_R^{(1)} = 0$,
which we have verified explicitly in~\eqref{eq:B2}.

This is the vector analogue of the Dirac result
$\beta_R^{(1/2)} = 0$ (from conformal invariance of the massless
Dirac action).  For the scalar, conformal invariance holds only
at $\xi = 1/6$, where $\beta_R^{(0)}(1/6) = 0$.

\subsection{Benchmark B10: no $\xi$-dependence}

Unlike the scalar field, the vector field has no free coupling
parameter.  The curvature coupling $\mathbf{U} = \mathrm{Ric}$
is uniquely fixed by gauge invariance (Weitzenbock identity).
Consequently, $h_C^{(1)}(x)$ and $h_R^{(1)}(x)$ are parameter-free
functions of~$x$ only.

\subsection{Benchmark B12: Gauss--Bonnet coefficient}

The $E_4$ coefficient of the Seeley--DeWitt $a_4$ for the physical
gauge field is
\begin{equation}
  \alpha_{E_4}^{(1)} = -\frac{31}{180}\,.
  \label{eq:B12}
\end{equation}
This is topological in $d = 4$ and does not contribute to
equations of motion.  It serves as a cross-check on the full
$a_4$ decomposition.

\subsection{Sign pattern and comparison across spins}

\begin{center}
\renewcommand{\arraystretch}{1.3}
\begin{tabular}{lccc}
\toprule
& \textbf{Scalar} ($s=0$, $\xi=0$)
& \textbf{Dirac} ($s=1/2$)
& \textbf{Vector} ($s=1$) \\
\midrule
$\tr(\Id)$ & $1$ & $4$ & $4$ \\
$\mathbf{U}$ & $\xi R$ & $R/4$ & $\mathrm{Ric}$ \\
$\Omega_{\mu\nu}$ & $0$ & spin curvature & Riemann \\
$\beta_W$ & $1/120$ & $-1/20$ & $1/10$ \\
$\beta_R$ & $1/72$ & $0$ & $0$ \\
$|\beta_W|/\beta_W^{(0)}$ & $1$ & $6$ & $12$ \\
UV: $h_C \sim$ & $+1/(12x)$ & $-1/(3x)$ & $-1/(3x)$ \\
Conformal? & at $\xi=1/6$ & yes & yes \\
\bottomrule
\end{tabular}
\end{center}

\begin{remark}[Ratio pattern]
The ratios $|\beta_W^{(s)}/\beta_W^{(0)}| = 1 : 6 : 12$ encode
both the fiber multiplicity ($\tr(\Id) = 1, 4, 4$) and the bundle
curvature enhancement.  For the scalar, $\Omega = 0$; for the
Dirac spinor, the spin connection gives a $3/2$ enhancement;
for the vector, the Riemann curvature gives a $3$-fold enhancement.
\end{remark}


%%======================================================================
\section{Spectral Action Form Factors $F_1^{(1)}$ and $F_2^{(1)}$}
\label{sec:spectral}
%%======================================================================

\subsection{From heat kernel to spectral action}

Following the NT-1 framework, the spectral action form factors are:
\begin{equation}
  F_i(z) = \frac{1}{16\pi^2}
  \int_0^\infty \dd t\,\tilde{f}(t)\,h_i(tz)\,,
  \qquad z = \Box/\Lambda^2\,,
  \label{eq:Fi-def}
\end{equation}
with the spectral function and its antiderivatives:
\begin{equation}
  \psi(u) = \int_0^\infty \dd t\,\tilde{f}(t)\,e^{-tu}\,,\quad
  \Psi_1(u) = \int_u^\infty \psi(v)\,\dd v\,,\quad
  \Psi_2(u) = \int_u^\infty (v-u)\,\psi(v)\,\dd v\,.
  \label{eq:psi-defs}
\end{equation}


\subsection{Spectral integration identities}

We use the identities from NT-1 (Lemma~3.2):
\begin{align}
  \int_0^\infty \dd t\,\tilde{f}(t)\,\varphi(tz)
  &= \int_0^1 \dd\alpha\,\psi(\alpha(1-\alpha)z)\,,
  \label{eq:int-phi}\\[4pt]
  \int_0^\infty \dd t\,\tilde{f}(t)\,\frac{1}{tz}
  &= \frac{\Psi_1(0)}{z}\,,
  \label{eq:int-1-over-x}\\[4pt]
  \int_0^\infty \dd t\,\tilde{f}(t)\,\frac{\varphi(tz)-1}{(tz)^2}
  &= \frac{1}{z^2}\int_0^1 \dd\alpha\,\bigl[
  \Psi_2(\alpha(1-\alpha)z) - \Psi_2(0)\bigr]\,.
  \label{eq:int-phim1-over-x2}
\end{align}

Additionally, we need:
\begin{equation}
  \int_0^\infty \dd t\,\tilde{f}(t)\,\frac{\varphi(tz)}{tz}
  = \frac{1}{z}\int_0^1 \dd\alpha\,\Psi_1(\alpha(1-\alpha)z)\,.
  \label{eq:int-phi-over-x}
\end{equation}


\subsection{$F_1^{(1)}$: the Weyl sector}

\begin{theorem}[Vector spectral form factor $F_1^{(1)}$]
\label{thm:F1}
\begin{equation}
  \boxed{
  \begin{aligned}
  F_1^{(1)}(z) &= \frac{1}{16\pi^2}\Biggl\{
    \frac{1}{4}\int_0^1\!\dd\alpha\,\psi(\alpha(1-\alpha)z)
    - \frac{5\,\Psi_1(0)}{6z}\\
  &\quad
    + \frac{1}{z}\int_0^1\!\dd\alpha\,\Psi_1(\alpha(1-\alpha)z)
    + \frac{1}{z^2}\int_0^1\!\dd\alpha\,
      \bigl[\Psi_2(\alpha(1-\alpha)z) - \Psi_2(0)\bigr]
  \Biggr\}
  \end{aligned}
  }
  \label{eq:F1}
\end{equation}
\end{theorem}

\begin{proof}
From $h_C^{(1)}(x) = \varphi/4 + (6\varphi-5)/(6x) + (\varphi-1)/x^2$,
we decompose each term:
\begin{align*}
  \frac{\varphi}{4} &\to \frac{1}{4}\int_0^1\psi(\alpha(1-\alpha)z)\,
  \dd\alpha\,,\\
  \frac{6\varphi-5}{6x} &= \frac{\varphi}{x} - \frac{5}{6x}
  \to \frac{1}{z}\int_0^1\Psi_1(\alpha(1-\alpha)z)\,\dd\alpha
  - \frac{5\Psi_1(0)}{6z}\,,\\
  \frac{\varphi-1}{x^2} &\to
  \frac{1}{z^2}\int_0^1[\Psi_2(\alpha(1-\alpha)z) - \Psi_2(0)]\,\dd\alpha\,.
\end{align*}
\end{proof}


\subsection{$F_2^{(1)}$: the $R^2$ sector}

\begin{theorem}[Vector spectral form factor $F_2^{(1)}$]
\label{thm:F2}
\begin{equation}
  \boxed{
  \begin{aligned}
  F_2^{(1)}(z) &= \frac{1}{16\pi^2}\Biggl\{
    -\frac{1}{48}\int_0^1\!\dd\alpha\,\psi(\alpha(1-\alpha)z)
    + \frac{11\,\Psi_1(0)}{72z}\\
  &\quad
    - \frac{1}{12z}\int_0^1\!\dd\alpha\,\Psi_1(\alpha(1-\alpha)z)
    + \frac{5}{12z^2}\int_0^1\!\dd\alpha\,
      \bigl[\Psi_2(\alpha(1-\alpha)z) - \Psi_2(0)\bigr]
  \Biggr\}
  \end{aligned}
  }
  \label{eq:F2}
\end{equation}
\end{theorem}

\begin{proof}
From $h_R^{(1)}(x) = -\varphi/48 + (11-6\varphi)/(72x)
+ 5(\varphi-1)/(12x^2)$:
\begin{align*}
  -\frac{\varphi}{48} &\to -\frac{1}{48}\int_0^1\psi\,\dd\alpha\,,\\
  \frac{11-6\varphi}{72x} &= -\frac{\varphi}{12x}
  + \frac{11}{72x}\\
  &\to -\frac{1}{12z}\int_0^1\Psi_1(\alpha(1-\alpha)z)\,\dd\alpha
  + \frac{11\Psi_1(0)}{72z}\,,\\
  \frac{5(\varphi-1)}{12x^2} &\to
  \frac{5}{12z^2}\int_0^1[\Psi_2(\alpha(1-\alpha)z) - \Psi_2(0)]\,\dd\alpha\,.
\end{align*}
\end{proof}


\subsection{Consistency check: exponential spectral function}

For $\psi(u) = e^{-u}$: $\Psi_1(u) = \Psi_2(u) = e^{-u}$,
$\Psi_1(0) = \Psi_2(0) = 1$.

All spectral integrals reduce to the $\varphi$-based form factors,
giving:
\begin{equation}
  F_i^{(1)}(z)\Big|_{\psi=e^{-u}}
  = \frac{h_i^{(1)}(z)}{16\pi^2}\,,\qquad i = 1, 2\,.
  \label{eq:exp-check}
\end{equation}

\verified{Numerical comparison of $F_i^{(1)}$ from the spectral
integral formulas against $h_i^{(1)}/(16\pi^2)$ for
$\psi = e^{-u}$ at 7 values of $z \in [0.01, 100]$.
All differences $< 10^{-14}$.}


%%======================================================================
\section{Limiting Cases}
\label{sec:limits}
%%======================================================================

\subsection{Flat space limit ($R = 0$)}

When $R_{\mu\nu\rho\sigma} = 0$, both $C^2 = 0$ and $R^2 = 0$,
so the quadratic curvature action~\eqref{eq:eff-action-weyl}
vanishes identically regardless of the form factors.
\quad$\checkmark$


\subsection{Massless limit}

The form factor argument $x = \Box/m^2$ is formal for a massless
gauge boson; in dimensional regularization the mass scale is replaced
by the renormalization scale $\mu$.  The closed forms remain valid
with $x \to -q^2/\mu^2$ for Euclidean momentum~$q$.


\subsection{Dimensional analysis}

In natural units ($c = \hbar = 1$):
\begin{itemize}[nosep]
\item The argument $x$ is dimensionless ($[\Box/m^2] = 1$).
\item $\varphi(x)$, $h_C^{(1)}(x)$, $h_R^{(1)}(x)$ are dimensionless.
\item $F_1^{(1)}(z)$ and $F_2^{(1)}(z)$ are dimensionless
  (checked: $[1/(16\pi^2)] \times [\psi] \times [\dd\alpha] = 1$).
\end{itemize}
\quad$\checkmark$


\subsection{Monotonicity and sign structure}

From numerical evaluation:
\begin{itemize}[nosep]
\item $h_C^{(1)}(x)$ is positive at $x = 0$ ($= 1/10$),
  decreases monotonically, crosses zero, and becomes negative,
  approaching $0^-$ as $x \to \infty$ (leading term $-1/(3x) < 0$).
\item $h_R^{(1)}(x)$ starts at $0$ ($x = 0$), increases to a
  positive maximum, then decreases toward $0^+$ as $x \to \infty$
  (leading term $+1/(9x) > 0$).
\end{itemize}

\verified{Sign structure confirmed numerically at 50 logarithmically
spaced points in $[10^{-3}, 10^6]$.}


%%======================================================================
\section{Summary of Results}
\label{sec:summary}
%%======================================================================

\subsection{Heat kernel form factors (boxed)}

\begin{tcolorbox}[colback=blue!5!white,colframe=blue!50!black,
  title={\textbf{Vector heat kernel form factors}}]
\textbf{Weyl sector} (parameter-free):
\begin{equation}
  h_C^{(1)}(x) = \frac{\varphi(x)}{4}
  + \frac{6\varphi(x) - 5}{6x}
  + \frac{\varphi(x) - 1}{x^2}\,,\qquad
  h_C^{(1)}(0) = \frac{1}{10}\,.
  \tag{\ref{eq:hC-final}}
\end{equation}

\textbf{$R^2$ sector} (parameter-free):
\begin{equation}
  h_R^{(1)}(x) = -\frac{\varphi(x)}{48}
  + \frac{11 - 6\varphi(x)}{72x}
  + \frac{5(\varphi(x)-1)}{12x^2}\,,\qquad
  h_R^{(1)}(0) = 0\,.
  \tag{\ref{eq:hR-final}}
\end{equation}

\textbf{Unconstrained (before ghost subtraction):}
\begin{align*}
  h_C^{(1,\mathrm{unconstr})} &= 2f_{\mathrm{Ric}}
  + \tfrac{1}{2}f_{\mathbf{U}} - 2f_\Omega\,,
  &h_C^{(1,\mathrm{unconstr})}(0) &= \frac{7}{60}\,.\\
  h_R^{(1,\mathrm{unconstr})} &= \tfrac{4}{3}f_{\mathrm{Ric}}
  + 4f_R + f_{R\mathbf{U}}
  + \tfrac{1}{3}f_{\mathbf{U}} - \tfrac{1}{3}f_\Omega\,,
  &h_R^{(1,\mathrm{unconstr})}(0) &= \frac{1}{36}\,.
\end{align*}

\textbf{Ghost subtraction:}
$h^{(1)} = h^{(1,\mathrm{unconstr})} - 2\,h^{(0)}(\xi=0)$
(two real FP ghosts).

\textbf{Master function:}
$\displaystyle\varphi(x) = \int_0^1 e^{-\alpha(1-\alpha)x}\,\dd\alpha$.
\end{tcolorbox}


\subsection{Spectral action form factors (boxed)}

\begin{tcolorbox}[colback=green!5!white,colframe=green!50!black,
  title={\textbf{Vector spectral action form factors}}]
\begin{equation}
  \begin{aligned}
  F_1^{(1)}(z) &= \frac{1}{16\pi^2}\Biggl\{
    \frac{1}{4}\int_0^1\!\psi(\alpha(1-\alpha)z)\,\dd\alpha
    - \frac{5\Psi_1(0)}{6z}\\
  &\quad
    + \frac{1}{z}\int_0^1\!\Psi_1(\alpha(1-\alpha)z)\,\dd\alpha
    + \frac{1}{z^2}\int_0^1\![\Psi_2(\alpha(1-\alpha)z) - \Psi_2(0)]
    \,\dd\alpha
  \Biggr\}
  \end{aligned}
  \tag{\ref{eq:F1}}
\end{equation}

\begin{equation}
  \begin{aligned}
  F_2^{(1)}(z) &= \frac{1}{16\pi^2}\Biggl\{
    -\frac{1}{48}\int_0^1\!\psi(\alpha(1-\alpha)z)\,\dd\alpha
    + \frac{11\Psi_1(0)}{72z}\\
  &\quad
    - \frac{1}{12z}\int_0^1\!\Psi_1(\alpha(1-\alpha)z)\,\dd\alpha
    + \frac{5}{12z^2}\int_0^1\![\Psi_2(\alpha(1-\alpha)z) - \Psi_2(0)]
    \,\dd\alpha
  \Biggr\}
  \end{aligned}
  \tag{\ref{eq:F2}}
\end{equation}

\textbf{Exponential check}: For $\psi = e^{-u}$:
$F_i^{(1)} = h_i^{(1)}/(16\pi^2)$.
\end{tcolorbox}


\subsection{Complete benchmark table}

\begin{center}
\renewcommand{\arraystretch}{1.3}
\begin{tabular}{llccl}
\toprule
\textbf{ID} & \textbf{Quantity} & \textbf{Target} & \textbf{Result} & \textbf{Reference} \\
\midrule
B1 & $\beta_W^{(1)} = h_C^{(1)}(0)$ & $1/10$ & $1/10$ & Eq.~\eqref{eq:B1} \\
B2 & $\beta_R^{(1)} = h_R^{(1)}(0)$ & $0$ & $0$ & Eq.~\eqref{eq:B2} \\
B3 & $h_C^{(1,\mathrm{unconstr})}(0)$ & $7/60$ & $7/60$ & Eq.~\eqref{eq:unconstr-limits} \\
B4 & $h_R^{(1,\mathrm{unconstr})}(0)$ & $1/36$ & $1/36$ & Eq.~\eqref{eq:unconstr-limits} \\
B5 & $h_C^{(1)} \to -1/(3x)$ & confirmed & confirmed & Eq.~\eqref{eq:hC-asymp} \\
B6 & $h_R^{(1)} \to 1/(9x)$ & confirmed & confirmed & Eq.~\eqref{eq:hR-asymp} \\
B7 & $h_C^{(1)} = 1/10 - 11x/420$ & confirmed & confirmed & Eq.~\eqref{eq:taylor-results} \\
B8 & $h_R^{(1)} = x/630 + O(x^2)$ & confirmed & confirmed & Eq.~\eqref{eq:taylor-results} \\
B9 & Ghost: $h^{(1)} - 2h^{(0)}(\xi\!=\!0)$ & confirmed & confirmed & Section~\ref{sec:ghost} \\
B10 & No $\xi$-dependence & confirmed & confirmed & Remark~\ref{rem:differences} \\
B11 & $\beta_R^{(1)}=0$ (conformal) & confirmed & confirmed & Section~\ref{sec:conformal} \\
B12 & $\alpha_{E_4}^{(1)} = -31/180$ & confirmed & confirmed & Eq.~\eqref{eq:B12} \\
\midrule
\multicolumn{5}{c}{Flat space $\checkmark$, Dimensionless $\checkmark$,
$F_i = h_i/(16\pi^2)$ for $\psi=e^{-u}$ $\checkmark$} \\
\bottomrule
\end{tabular}
\end{center}


%%======================================================================
\begin{thebibliography}{20}

\bibitem{CZ2012}
A.~Codello and O.~Zanusso,
``On the non-local heat kernel expansion,''
J.\ Math.\ Phys.\ \textbf{54} (2013) 013513
[arXiv:1203.2034 [hep-th]].

\bibitem{TSR2020}
P.~de~Morais~Teixeira, I.~L.~Shapiro, and T.~G.~Ribeiro,
``One-loop effective action: nonlocal form factors and
renormalization group,''
Grav.\ Cosmol.\ \textbf{26} (2020) 185
[arXiv:2003.04503 [hep-th]].

\bibitem{Vassilevich2003}
D.~V.~Vassilevich,
``Heat kernel expansion: user's manual,''
Phys.\ Rept.\ \textbf{388} (2003) 279
[hep-th/0306138].

\bibitem{BV1990a}
A.~O.~Barvinsky and G.~A.~Vilkovisky,
``Covariant perturbation theory (II).\ Second order in the curvature.
General algorithms,''
Nucl.\ Phys.\ B \textbf{333} (1990) 471.

\bibitem{CPR2009}
A.~Codello, R.~Percacci, and C.~Rahmede,
``Investigating the ultraviolet properties of gravity with a
Wilsonian renormalization group equation,''
Ann.\ Phys.\ \textbf{324} (2009) 414
[arXiv:0805.2909 [hep-th]].

\bibitem{NT1b_vector_conventions}
D.~Alfyorov,
``NT-1b Vector (Gauge Boson) Form Factors: Literature Conventions
and Benchmarks,''
working document, March 2026.

\bibitem{NT1b_scalar}
D.~Alfyorov,
``NT-1b Scalar Form Factors: Main Derivation,''
working document, March 2026.

\bibitem{NT1}
D.~Alfyorov,
``NT-1: Nonlocal Form Factors of the SCT Spectral Action ---
From Heat Kernel to Spectral Action for the Dirac Operator,''
working document, March 2026.

\end{thebibliography}

\end{document}
